\documentclass[11pt,a4paper]{article}
\usepackage[utf8]{inputenc}
\usepackage[T1]{fontenc}
\usepackage{amsmath,amssymb}
\usepackage{hyperref}

\title{Atlas: A Line-Oriented DSL for Prototyping Historical Risk Metrics}
\author{Atlas contributors}
\date{\today}

\begin{document}
\maketitle

\begin{abstract}
We describe Atlas, a minimal domain-specific language for scripting
historical return transforms, rolling volatility, value-at-risk (VaR),
and expected shortfall (ES) on columnar market data. Atlas emphasizes
a small assignment-oriented surface syntax, a strict separation between
parsing and statistical primitives, and an extensible standard library
layout suitable for teaching and rapid experimentation.
\end{abstract}

\section{Introduction}
Market risk workflows often repeat the same sequence of operations:
ingest tabular prices, construct returns, estimate dispersion and tail
risk, and report scalar or vector summaries. General-purpose notebooks
are flexible but can obscure assumptions; Atlas proposes a compact
intermediate representation that keeps dataflow explicit.

\section{Language surface}
Programs are newline-separated statements. The primary forms are
assignment (\texttt{name = func primary \ldots}) and printing
(\texttt{print name}). Relative file paths in \texttt{load} resolve
from the script directory to make examples portable.

\section{Estimators}
Let $\{r_t\}$ denote a return series. Rolling volatility uses the
sample standard deviation over a sliding window. Historical VaR at
confidence level $\alpha$ reports the empirical $(1-\alpha)$ quantile.
Historical ES reports the sample mean of observations at or below the
VaR threshold.

\section{Implementation}
The reference implementation is in Python: parsing produces an
immutable command record; an interpreter maintains a name-to-object
environment mapping strings to \texttt{pandas.DataFrame} values; and
numeric kernels live in a dedicated standard-library module tree.

\section{Conclusion}
Atlas is intentionally small. Future work includes parametric and
simulation-based VaR, portfolio-level aggregation, and optional static
checks for schema and units.

\end{document}
